% ============================================================
%  GPF Plume Navigation — Experimental Setup
%  Drop-in section for the Experiments chapter.
%  Assumes GPF framework has already been introduced.
% ============================================================

\subsection{Plume Navigation Task}

We evaluate the GPF framework on chemical plume navigation, the problem of
locating an odor source by following a turbulent plume back to its origin.
This task is representative of the class of sparse-reward, partially
observable navigation problems that motivate adaptive-depth networks: the
signal is intermittent, the environment is non-stationary, and the
appropriate behavioral strategy (surge, cast, or loiter) depends on
context that must be inferred from recent sensory history.

\subsubsection{Filament-Based Plume Simulator}

We simulate the plume using a Farrell--Murlis filament
model~\cite{farrell2002filament}.  The environment is a
$20\,\text{m} \times 20\,\text{m}$ two-dimensional domain operating at a
control frequency of $10\,\text{Hz}$ ($\Delta t = 0.1\,\text{s}$).

\paragraph{Source and filament emission.}
At the start of each episode the source is placed on the upwind half of the
domain at a randomly sampled position biased toward the upwind edge
($\pm 25\%$ of domain width from center, with an additional $\pm 15\%$
jitter).  The source emits discrete filaments at a mean rate of
$\lambda = 5\,\text{filaments/s}$, drawn from a Poisson process so that
$n_{\text{new}} \sim \text{Poisson}(\lambda\,\Delta t)$ filaments are
released each time step.  Each filament is initialized at the source with a
small positional jitter drawn from $\mathcal{N}(0,\,\sigma_0^2)$ with
$\sigma_0 = 0.01\,\text{m}$.

\paragraph{Filament dynamics.}
Filaments advect with the local wind velocity $(v_x, v_y)$ and grow by
molecular diffusion.  The Gaussian radius of filament $i$ at age
$\tau_i$ is
\begin{equation}
  \sigma_i(\tau_i) = \sqrt{\sigma_0^2 + 2\,D\,\tau_i},
  \label{eq:filament_diffusion}
\end{equation}
with diffusion rate $D = 0.05\,\text{m}^2/\text{s}$.  Filaments are
removed from the simulation after a decay time of $30\,\text{s}$ or when
they leave the domain boundary.

\paragraph{Concentration measurement.}
The concentration contributed by filament $i$ at sensor position
$\mathbf{p}$ is
\begin{equation}
  c_i(\mathbf{p}) = \frac{m}{2\pi\sigma_i^2}
    \exp\!\left(-\frac{\|\mathbf{p} - \mathbf{x}_i\|^2}{2\sigma_i^2}\right),
  \label{eq:filament_conc}
\end{equation}
where $m = 0.1$ is the filament mass and $\mathbf{x}_i$ is the filament
center.  Filaments beyond $3\sigma_i$ are excluded for computational
efficiency.  The total concentration is the sum of all active filament
contributions, clamped to a saturation ceiling of $1.0$ and corrupted by
additive Gaussian noise $\mathcal{N}(0,\,\sigma_n^2)$ with
$\sigma_n = 10^{-3}$.

\paragraph{Wind model.}
Wind speed $v$ and direction $\theta$ follow coupled discrete-time
Ornstein--Uhlenbeck processes with mean-reversion coefficient
$\alpha = \Delta t / \tau_c$:
\begin{align}
  v_{t+1}   &= v_t - \alpha(v_t - \bar{v}) +
               \sigma_v \sqrt{2\alpha}\;\epsilon_v, \\
  \theta_{t+1} &= \theta_t - \alpha(\theta_t - \bar{\theta}) +
               \sigma_\theta \sqrt{2\alpha}\;\epsilon_\theta,
\end{align}
where $\epsilon_v, \epsilon_\theta \sim \mathcal{N}(0,1)$ are independent,
$\bar{v} = 1.0\,\text{m/s}$, $\bar{\theta} = 0\,\text{rad}$,
$\sigma_v = 0.2\,\text{m/s}$, $\sigma_\theta = 15^\circ$, and correlation
time $\tau_c = 2\,\text{s}$.  Wind speed is clamped to a minimum of
$0.1\,\text{m/s}$.

\subsubsection{Agent and Sensor Model}

The agent carries a bilateral antenna pair modeled as two point sensors in
the body frame.  The left and right antennae are offset $0.04\,\text{m}$
laterally from the centerline and $0.08\,\text{m}$ forward of the agent
origin, giving an inter-antenna separation of $0.1\,\text{m}$.

At each time step the agent selects one of six discrete actions:
\textit{forward}, \textit{turn left} $15^\circ$, \textit{turn right}
$15^\circ$, \textit{turn around} $180^\circ$, \textit{cast left}
$30^\circ$, and \textit{cast right} $30^\circ$.  All actions are followed
by a forward step of $v = 0.5\,\text{m/s} \times \Delta t = 0.05\,\text{m}$.
Domain boundaries are reflecting: the agent position is clipped to
$[0, 20\,\text{m}]^2$ rather than terminating the episode, providing
continuous gradient signal and a non-trivial reactive baseline.  An episode
terminates upon success (agent within $r_s = 0.5\,\text{m}$ of the source)
or after a budget of $T_{\max} = 1200$ steps ($100\,\text{s}$).

At the start of each episode the agent is placed $3$--$10\,\text{m}$
downwind of the source, oriented toward the source with heading noise
$\mathcal{N}(0,\,(30^\circ)^2)$.

\subsubsection{State Representation}

Raw sensor observations are discretized into a compact state tuple
$(b_L,\, b_R,\, w)$ where $b_L, b_R \in \{0,\ldots,6\}$ are concentration
bin indices for the left and right antennae, and $w \in \{0,\ldots,7\}$ is
a wind-direction octant relative to the agent heading.

Concentration is quantized into seven levels: a \textsc{blank} level for
readings at or below the noise floor, and six log-spaced bins with absolute
edges at $[3, 6, 15, 45, 150, 600] \times \sigma_n$
$= [0.003, 0.006, 0.015, 0.045, 0.15, 0.6]$.  Wind octant $0$ corresponds
to a headwind (wind arriving from directly ahead); octants increase
clockwise.  Calm conditions ($v < 0.05\,\text{m/s}$) default to octant $0$.

The state tuple is one-hot encoded into a $22$-dimensional binary vector
($7 + 7 + 8$) that serves as the Q-network input.  The discrete state space
has $7 \times 7 \times 8 = 392$ unique states.

\subsubsection{GPF Expected SARSA Agent}

The Q-function is represented by an MLP with ReLU activations.  The initial
architecture is $22 \to 64 \to 6$ (one hidden layer of width $64$).  The
network is trained with the Expected SARSA update rule under an
$\varepsilon$-greedy policy, with $\varepsilon$ decayed from $1.0$ to
$0.05$ multiplicatively by a factor of $0.9995$ per episode.  The optimizer
is Adam with learning rate $\eta = 10^{-3}$, and the discount factor is
$\gamma = 0.99$.

The GPF grow trigger monitors a split-window EMA of the evaluation success
rate: a grow event fires when the improvement from the first half to the
second half of the EMA history falls below $\varepsilon_{\text{add}} = 0.05$
(5\%), subject to a warm-up of $400$ episodes and a per-grow cooldown of
$400$ episodes.  New layers are initialized near-identity
($W = I + 0.01\,\mathcal{N}$) to minimize policy disruption.
The maximum network depth is four hidden layers.

Pruning uses Optimal Brain Damage
saliency~\cite{lecun1990optimal}: a weight $w_{ij}$ in layer $l$ is zeroed
if $\tfrac{1}{2}w_{ij}^2\,\mathbb{E}[g_{ij}^2] < \tau_{\text{prune}}$,
where $\mathbb{E}[g_{ij}^2]$ is estimated from $1500$ accumulated gradient
steps and $\tau_{\text{prune}} = 10^{-6}$.  Pruning fires once per grow
event on the preceding layers.  Freeze is configured but disabled
($\text{min\_episodes\_before\_freeze} = 99{,}999$) so that all performance
is attributable to grow and prune alone.

\subsubsection{Reward Function}

The reward at each step is
\begin{equation}
  r_t = r_{\text{time}} + r_{\text{event}} + r_{\text{shape}},
\end{equation}
where:
\begin{itemize}
  \item $r_{\text{time}} = -0.01$ per step (time penalty encouraging efficiency),
  \item $r_{\text{event}} = +100$ on source discovery; $+1.0$ if either
    antenna detects a whiff (concentration $> 3\sigma_n$); $-0.05$ if both
    antennae have been blank for more than $2\,\text{s}$ ($20$ steps),
  \item $r_{\text{shape}} = 2.0\,(d_{t-1} - d_t)$ is a potential-based
    distance-shaping term~\cite{ng1999policy} that provides dense gradient
    signal toward the source without altering the optimal policy.
\end{itemize}

\subsubsection{Training Protocol}

All experiments use seed $42$ for the environment and seed $43$ for the
agent.  Training runs for $5{,}000$ episodes with evaluation every $500$
episodes over $200$ held-out episodes.  The best checkpoint (highest
evaluation success rate) is retained.  Key hyperparameters are summarized
in Table~\ref{tab:gpf_hyperparams}.

\begin{table}[h]
\centering
\caption{Hyperparameters for the plume navigation GPF agent.}
\label{tab:gpf_hyperparams}
\begin{tabular}{lll}
\toprule
\textbf{Parameter} & \textbf{Symbol} & \textbf{Value} \\
\midrule
\multicolumn{3}{l}{\textit{Environment}} \\
Domain size           & ---                    & $20 \times 20\,\text{m}$ \\
Time step             & $\Delta t$             & $0.1\,\text{s}$ \\
Episode budget        & $T_{\max}$             & $1200$ steps ($100\,\text{s}$) \\
Source emission rate  & $\lambda$              & $5\,\text{filaments/s}$ \\
Success radius        & $r_s$                  & $0.5\,\text{m}$ \\
Mean wind speed       & $\bar{v}$              & $1.0\,\text{m/s}$ \\
Wind correlation time & $\tau_c$               & $2.0\,\text{s}$ \\
\midrule
\multicolumn{3}{l}{\textit{Agent / network}} \\
Hidden width          & $d_h$                  & $64$ \\
Max hidden layers     & $L_{\max}$             & $4$ \\
Learning rate         & $\eta$                 & $10^{-3}$ \\
Discount factor       & $\gamma$               & $0.99$ \\
Initial exploration   & $\varepsilon_0$        & $1.0$ \\
Final exploration     & $\varepsilon_\infty$   & $0.05$ \\
Exploration decay     & ---                    & $0.9995$ per episode \\
\midrule
\multicolumn{3}{l}{\textit{GPF}} \\
Grow threshold        & $\varepsilon_{\text{add}}$ & $0.05$ \\
Grow warm-up / cooldown & $\varepsilon_e / \varepsilon_k$ & $400$ episodes each \\
Prune threshold       & $\tau_{\text{prune}}$  & $10^{-6}$ \\
Prune accumulation    & ---                    & $1500$ steps \\
\bottomrule
\end{tabular}
\end{table}
